\section{Process composition and return values}

\subsection{\code{MAIN-PROCESS} (\implemented)}

\begin{lstlisting}
MAIN-PROCESS FullAdder
\end{lstlisting}
This marks the process body and supplies its catalog/library name. Compilation
already begins from the named process, so the row itself emits no gate.
The name is the first positional word after \code{MAIN-PROCESS}; avoid spaces
inside it. This name connects a protocol to catalog entries, child calls,
downloads, and companion truth-table metadata. If no marker exists, the system
uses \code{InlineProcess}, but named processes are strongly recommended.

\subsection{\code{DECLARECHILD} (\acceptedonly)}

\begin{lstlisting}
DECLARECHILD SingleBitFullAdder
\end{lstlisting}
The declaration is logged. It does not load inline source. Child bodies must
exist in the external \code{librarySources} catalog when called.
DECLARECHILD is therefore documentation/compatibility syntax rather than an
import statement. A successful declaration does not guarantee that a later
call can resolve. The process catalog supplied to compilation remains the
source of truth.

\subsection{\code{RUNCHILD} --- expand a child process (\implemented)}

\begin{lstlisting}
RUNCHILD SingleBitFullAdder \
  -I A B Cin \
  -O Carry Sum
\end{lstlisting}
The first positional argument is the child process name. Inputs bind by the
child's PARAMS order: A maps to the first parameter, B to the second, and Cin
to the third in this example. Outputs bind by the child's RETURNVALS order.
Before expansion, each bound parent output is zero-prepared once. Child-local
names receive a unique scope.

RUNCHILD is compile-time composition, not a dynamic subroutine call in a
running quantum machine. The child body is recursively expanded into the
parent's flat gate sequence. This allows the renderer and simulator to inspect
every resulting gate.

\subsection{\code{CALL} --- alternate child invocation (\implemented)}

\begin{lstlisting}
CALL SingleBitFullAdder -I A B Cin -O Carry Sum
\end{lstlisting}
CALL currently follows the same child lookup, parameter binding, output
preparation, and expansion path as RUNCHILD. RUNCHILD is used by the bundled
multi-bit protocols and makes compile-time expansion explicit; CALL is retained
as an equivalent system spelling. Both fail when the named child is absent from
the supplied process library.

\subsection{\code{RETURNVALS} (\implemented)}

\begin{lstlisting}
RETURNVALS Carry Sum
\end{lstlisting}
Positional tokens define the process's ordered public outputs. They determine
child output binding, truth-table output columns, and the UI's preferred
logical width. Do not use \code{-I} or \code{-O} here: flags would be retained
as positional return tokens.

RETURNVALS does not measure automatically. Add MEASURE instructions when
classical outcomes are required. Returning an unmeasured wire preserves it as
part of the projected quantum output shown by the system.

\subsection{\code{ACCEPTVALS} (\implemented)}

\begin{lstlisting}
CALL Child -I A B
ACCEPTVALS LocalCarry LocalSum
\end{lstlisting}
Each positional local token aliases the corresponding value from the most
recent child return list. Direct \code{-O} binding on RUNCHILD/CALL is clearer
when output names are known.

Ordering is again essential: the first ACCEPTVALS name receives the first
returned reference, and so on. The instruction refers to the most recently
completed child expansion in compiler state; inserting another child call
changes which return list is accepted.

\subsection{Nested full-adder example}

\begin{lstlisting}
# child.qpucir
PARAMS: A:state B:state Cin:state
MAIN-PROCESS ChildAdder
CREATETOKEN -I Carry Sum
SET Carry 0p
SET Sum 0p
XOR   -I A B   -O Sum
CNOT -I Cin   -O Sum
CCNOT -I A B   -O Carry
CCNOT -I A Cin -O Carry
CCNOT -I B Cin -O Carry
RETURNVALS Carry Sum
\end{lstlisting}

\begin{lstlisting}
# parent.qpucir
PARAMS: A:state B:state Cin:state
MAIN-PROCESS Parent
CREATETOKEN -I Cout Result
RUNCHILD ChildAdder -I A B Cin -O Cout Result
MEASURE -I Cout
MEASURE -I Result
RETURNVALS Cout Result
\end{lstlisting}

For basis inputs, the full-adder equations are
\[
\mathrm{Sum}=A\oplus B\oplus C_{\mathrm{in}},\qquad
C_{\mathrm{out}}=(A\land B)\oplus(A\land C_{\mathrm{in}})
\oplus(B\land C_{\mathrm{in}}).
\]
\begin{center}
\begin{tabular}{ccc|cc}\toprule
$A$&$B$&$C_{\rm in}$&$C_{\rm out}$&Sum\\\midrule
0&0&0&0&0\\0&0&1&0&1\\0&1&0&0&1\\0&1&1&1&0\\
1&0&0&0&1\\1&0&1&1&0\\1&1&0&1&0\\1&1&1&1&1\\\bottomrule
\end{tabular}
\end{center}

